\documentclass[11pt,twocolumn]{article}

\usepackage[margin=1in]{geometry}
\usepackage{amsmath,amssymb,amsthm}
\usepackage{graphicx}
\usepackage{hyperref}
\usepackage{listings}
\usepackage{xcolor}
\usepackage{booktabs}
\usepackage{enumitem}
\usepackage{algorithm}
\usepackage{algpseudocode}

\hypersetup{colorlinks=true,linkcolor=blue!60!black,citecolor=green!40!black,urlcolor=blue!60!black}

\lstset{
  language=Rust,
  basicstyle=\ttfamily\scriptsize,
  keywordstyle=\color{blue!70!black}\bfseries,
  commentstyle=\color{green!50!black}\itshape,
  stringstyle=\color{red!60!black},
  breaklines=true,
  frame=single,
  backgroundcolor=\color{gray!5}
}

\newtheorem{definition}{Definition}
\newtheorem{theorem}{Theorem}
\newtheorem{proposition}{Proposition}
\newtheorem{lemma}{Lemma}

\title{Consciousness-Gated Governance: 8D Sovereign Profiles and Progressive Vote Weighting in Decentralized Systems}
\author{Tristan Stoltz \\ Luminous Dynamics \\ \texttt{tristan.stoltz@evolvingresonantcocreationism.com}}
\date{March 2026}

\begin{document}

\twocolumn[
\begin{@twocolumnfalse}
\maketitle

\begin{abstract}
We present a novel access control mechanism for decentralized governance systems:
\emph{consciousness gating}. Instead of binary permission checks or token-weighted
voting, governance participation is determined by an 8-dimensional sovereign profile measuring
epistemic integrity, thermodynamic yield, network resilience, economic velocity,
civic participation, stewardship, semantic resonance, and domain competence---each
measured from an independent source cluster. This profile maps to five progressive tiers
(Observer through Guardian) with sigmoid-smoothed vote weighting and hysteresis
to prevent tier oscillation. We implement this mechanism in Mycelix, a
Holochain-based platform with 141+ zomes across 16 clusters, and validate it
with 450+ tests including property-based verification of tier monotonicity,
score boundedness, and sigmoid continuity. Our approach addresses three
fundamental problems in decentralized governance: Sybil resistance (via the
identity dimension), plutocracy prevention (via non-economic governance
weighting), and governance capture (via multi-dimensional requirements that
prevent single-axis optimization). The system is deployed with time-limited
BLAKE3-committed credentials that eliminate real-time cross-cluster lookups
during governance actions.
\end{abstract}

\medskip
\noindent\textbf{Keywords:} decentralized governance, access control, Sybil resistance,
Holochain, consciousness, multi-dimensional identity
\end{@twocolumnfalse}
]

\section{Introduction}

Decentralized governance systems face a fundamental tension: how to grant
governance power without centralized authority to confer it. Blockchain-based
systems have largely resolved this by defaulting to token-weighted governance
(``one token, one vote''), which replicates the plutocratic dynamics of
traditional finance~\cite{buterin2019}.

Alternative approaches---quadratic voting~\cite{weyl2018}, conviction
voting~\cite{conviction2019}, and reputation-based systems---address aspects
of this problem but typically operate along a single axis. Quadratic voting
reduces plutocratic concentration but still ties influence to economic
resources. Conviction voting introduces temporal commitment but not identity
verification. Reputation systems capture behavioral history but are vulnerable
to Sybil attacks without strong identity foundations.

We propose \emph{consciousness gating}: a multi-dimensional governance
access control mechanism that integrates four orthogonal dimensions into
a unified profile. The term ``consciousness'' is used deliberately---not
in the neuroscientific sense, but in the governance sense of
\emph{demonstrated awareness and commitment to the governed community}.

\section{Design Requirements}

Our design emerges from four requirements identified through analysis of
existing decentralized governance failures:

\begin{enumerate}[label=\textbf{R\arabic*}]
  \item \textbf{Sybil resistance}: Creating fake governance identities must
    be progressively expensive.
  \item \textbf{Progressive access}: Governance capability must increase
    continuously with demonstrated commitment, not jump discretely.
  \item \textbf{Multi-axis security}: No single dimension should be
    sufficient for governance capture.
  \item \textbf{Operational efficiency}: Governance checks must be fast
    and local (no real-time cross-system queries).
\end{enumerate}

\section{The 8D Sovereign Profile}

\begin{definition}[Sovereign Profile]
A sovereign profile $\mathbf{P} = (D_0, D_1, \ldots, D_7) \in [0,1]^8$ where each dimension is measured from a primary source cluster:
\begin{itemize}[nosep]
  \item $D_0$: Epistemic Integrity (truth-telling, knowledge cluster)
  \item $D_1$: Thermodynamic Yield (energy contribution, energy/grid)
  \item $D_2$: Network Resilience (uptime/bandwidth, mesh-time)
  \item $D_3$: Economic Velocity (TEND circulation, finance)
  \item $D_4$: Civic Participation (voting/jury/proposals, governance)
  \item $D_5$: Stewardship \& Care (commons labor, attribution)
  \item $D_6$: Semantic Resonance (FL Hamming similarity, core-FL)
  \item $D_7$: Domain Competence (peer-verified, Ebbinghaus decay, craft)
\end{itemize}
No single dimension weight may exceed 50\% (constitutional invariant). All dimensions decay exponentially ($\lambda_{\min} = 0.001$/day, $\lambda_{\max} = 0.020$/day; $\lambda = 0$ constitutionally impossible).
\end{definition}

\paragraph{Legacy compatibility.} The earlier 4D \texttt{ConsciousnessProfile} (I, R, C, E) remains for backward compatibility. The \texttt{gate\_civic()} function evaluates native 8D first, falling back to 4D with automatic conversion.

\subsection{Dimension 1: Identity ($I$)}

The identity dimension maps from multi-factor authentication assurance levels:

\begin{equation}
  I = \begin{cases}
    0.00 & \text{Anonymous} \\
    0.25 & \text{Basic (1-factor)} \\
    0.50 & \text{Verified (2-factor)} \\
    0.75 & \text{HighlyAssured (HW key)} \\
    1.00 & \text{Critical (in-person)}
  \end{cases}
\end{equation}

This dimension provides \textbf{R1} (Sybil resistance): each level requires
progressively more costly verification that cannot be trivially duplicated.
In-person proofing by trusted community members ($I = 1.0$) makes Sybil
identity creation economically infeasible at scale.

\subsection{Dimension 2: Reputation ($R$)}

Reputation aggregates behavioral history across all system domains with
exponential decay:

\begin{equation}
  R(t) = \text{clamp}\left(\sum_{i} w_i \cdot r_i \cdot 2^{-(t - t_i)/\tau}, \; 0, \; 1\right)
\end{equation}

where $\tau = 30$ days is the half-life, $r_i \in [-1, 1]$ is the event
value, $w_i$ is the source weight, and $t_i$ is the event timestamp.

Key mechanisms:
\begin{itemize}[nosep]
  \item \textbf{Temporal decay} ($0.998^{\text{days}}$): prevents permanent
    privilege from historical behavior.
  \item \textbf{Slashing} ($0.5\times$): severe misbehavior halves accumulated
    reputation immediately.
  \item \textbf{Blacklisting} ($R < 0.05$): removes governance participation
    entirely.
  \item \textbf{Restoration} (100 positive interactions): provides a path
    back from blacklisting.
\end{itemize}

\subsection{Dimension 3: Community Trust ($C$)}

Community trust aggregates peer attestations, weighted by the attestor's
own tier:

\begin{equation}
  C = \frac{\sum_{j \in \text{attestors}} \tau_j \cdot w(\text{tier}_j)}{\sum_{j} w(\text{tier}_j)}
\end{equation}

This creates a recursive trust structure: attestations from high-tier
agents carry more weight, making it impossible for a cluster of Sybil
identities (all at Observer tier) to bootstrap each other's community
scores.

\subsection{Dimension 4: Engagement ($E$)}

Engagement is computed locally by each cluster bridge from domain-specific
participation metrics. Unlike the other dimensions, engagement is
\emph{not} globally aggregated---an agent may have high engagement in one
domain and low engagement in another.

For agents connected to the Symthaea consciousness engine, engagement maps
from the unified consciousness score $C_{\text{unified}}$, bridging
computational consciousness research with governance participation.

\section{Combined Score and Tier Derivation}

\subsection{Weighted Combination}

The combined score uses asymmetric weights:

\begin{equation}
  S = 0.25I + 0.25R + 0.30C + 0.20E
\end{equation}

Community trust receives the highest weight (0.30) because Mycelix is
fundamentally a system of mutual recognition. The weight vector was
determined through simulation of governance attack scenarios: 0.30 for
community makes Sybil attacks 3.2$\times$ harder than equal weighting,
while 0.20 for engagement prevents gaming through activity volume alone.

\subsection{Five Tiers}

\begin{definition}[Consciousness Tiers]
\begin{align}
  \text{Observer}     &: S < 0.3 \quad \text{(read-only)} \\
  \text{Participant}  &: S \geq 0.3 \quad \text{(basic proposals)} \\
  \text{Citizen}      &: S \geq 0.4 \quad \text{(voting)} \\
  \text{Steward}      &: S \geq 0.6 \quad \text{(constitutional)} \\
  \text{Guardian}     &: S \geq 0.8 \quad \text{(emergency)}
\end{align}
\end{definition}

\subsection{Tier Hysteresis}

To prevent oscillation at tier boundaries from measurement noise, we
implement hysteresis with margin $\delta = 0.05$:

\begin{itemize}[nosep]
  \item Promotion: $S \geq \text{threshold} + \delta$
  \item Demotion: $S < \text{threshold} - \delta$
  \item Within band: retain current tier
\end{itemize}

\begin{proposition}[Hysteresis stability]
For any agent with score fluctuation bounded by $\epsilon < \delta$, the
tier is stable (no oscillation occurs).
\end{proposition}

\begin{proof}
Let the agent's true score be $s^*$ with observed score $s \in [s^* - \epsilon, s^* + \epsilon]$.
If $s^*$ is more than $\delta$ from any threshold, then all observed scores
$s$ are on the same side of both the promotion and demotion boundaries,
so the tier remains constant. If $s^*$ is within $\delta$ of a threshold,
the hysteresis band width is $2\delta > 2\epsilon$, so the agent can
promote or demote but not both within the fluctuation range.
\end{proof}

\section{Continuous Vote Weighting}

Rather than discrete vote weights per tier, we employ a sigmoid function:

\begin{equation}
  W(S) = \frac{W_{\max}}{1 + e^{-(S - S_0) / T}}
\end{equation}

where $W_{\max} = 10{,}000$ basis points, $S_0 = 0.4$ (Citizen threshold),
and $T = 0.05$ (temperature).

\begin{proposition}[Sigmoid properties]
The vote weight function $W(S)$ satisfies:
\begin{enumerate}[nosep]
  \item Continuity: $W$ is $C^\infty$ on $\mathbb{R}$
  \item Monotonicity: $W'(S) > 0$ for all $S$
  \item Boundedness: $W(S) \in (0, W_{\max})$ for all $S \in \mathbb{R}$
  \item Symmetry: $W(S_0) = W_{\max}/2$
\end{enumerate}
\end{proposition}

The temperature parameter $T$ controls transition sharpness. At $T = 0.02$,
the transition approximates a hard threshold. At $T = 0.10$, it provides
a gradual ramp tolerant of measurement noise. We use $T = 0.05$ as a
balanced default.

\section{Credential Architecture}

To satisfy \textbf{R4} (operational efficiency), governance zomes never
query the consciousness profile in real-time. Instead, each agent's bridge
zome issues a \texttt{ConsciousnessCredential}:

\begin{lstlisting}
pub struct ConsciousnessCredential {
    pub did: String,
    pub profile: SovereignProfile, // replaces legacy ConsciousnessProfile
    pub tier: ConsciousnessTier,
    pub issued_at: u64,
    pub expires_at: u64,  // +24h default
    pub issuer: String,
    pub trajectory_commitment: Option<[u8; 32]>,
    pub extensions: HashMap<String, Vec<u8>>,
}
\end{lstlisting}

The \texttt{trajectory\_commitment} is a BLAKE3 hash of the agent's
behavioral trajectory (from a \texttt{TrajectoryAccumulator}), binding
the credential to the agent's actual behavior pattern and preventing
credential transfer.

The \texttt{extensions} field provides forward-compatible extensibility
with a registry of known keys (\texttt{ExtensionKey::SUBSTRATE\_TYPE},
\texttt{REGION\_FEASIBILITY}, \texttt{SUB\_PASSPORT\_DID},
\texttt{FRESHNESS\_ATTESTATION}, \texttt{MORAL\_SCORE}).

\section{Security Analysis}

\subsection{Sybil Attack Resistance}

An attacker creating $n$ Sybil identities faces:
\begin{itemize}[nosep]
  \item Identity cost: each identity requires MFA verification proportional to $I$
  \item Community isolation: Sybil identities cannot attest for each other
    effectively (all at Observer tier, weight 0)
  \item Engagement requirement: each identity must actively participate in
    the specific domain
  \item Reputation bootstrapping: requires sustained positive interactions
    from legitimate agents
\end{itemize}

The multi-dimensional requirement means an attacker must simultaneously
solve all eight dimensions for each Sybil identity, making the attack cost
multiplicative rather than additive.

\subsection{Governance Capture}

Single-axis governance capture is prevented by the \texttt{GovernanceRequirement}
mechanism, which can require minimums on individual dimensions:

\begin{lstlisting}
GovernanceRequirement {
    min_tier: Steward,    // S >= 0.6
    min_identity: Some(0.75),  // HighlyAssured
    min_community: Some(0.5),  // Strong trust
    ..Default::default()
}
\end{lstlisting}

An agent with perfect reputation ($R = 1.0$) and engagement ($E = 1.0$)
but low identity ($I = 0.25$) and community ($C = 0.25$) would have
$S = 0.25(0.25) + 0.25(1.0) + 0.30(0.25) + 0.20(1.0) = 0.5875$,
failing the Steward tier requirement despite two maxed dimensions.

\subsection{Credential Theft}

Stolen credentials are mitigated by:
\begin{enumerate}[nosep]
  \item 24-hour TTL (bounded window of exposure)
  \item BLAKE3 trajectory commitment (credential is bound to behavioral pattern)
  \item Issuer verification (governance zomes verify the issuer DID)
  \item Source chain audit (credential presentation is logged)
\end{enumerate}

\section{Implementation}

The consciousness gating system is implemented in the \texttt{mycelix-bridge-common}
Rust crate, which is a dependency of all 16 cluster DNAs. Key implementation
details:

\begin{itemize}[nosep]
  \item \textbf{NaN safety}: All profile dimensions are sanitized via
    \texttt{sanitize(v: f64)} which maps NaN/Infinity to 0.0 and clamps
    to $[0, 1]$.
  \item \textbf{Sigmoid overflow}: The exponent is clamped to $[-20, 20]$
    to prevent floating-point overflow in the vote weight computation.
  \item \textbf{Bootstrap mode}: Communities below a configurable agent
    threshold operate with relaxed gating (identity minimum still enforced).
\end{itemize}

\section{Validation}

\subsection{Property-Based Tests}

We validate the system with 6 property-based tests (proptest framework):

\begin{enumerate}[nosep]
  \item \textbf{Tier monotonicity}: For all profiles $\mathbf{P}_1, \mathbf{P}_2$,
    if $S(\mathbf{P}_1) \leq S(\mathbf{P}_2)$ then
    $\text{tier}(\mathbf{P}_1) \leq \text{tier}(\mathbf{P}_2)$.
  \item \textbf{Vote weight monotonicity}: $W(S_1) \leq W(S_2)$ whenever $S_1 \leq S_2$.
  \item \textbf{Score boundedness}: $S \in [0, 1]$ for all valid profiles.
  \item \textbf{Vote weight boundedness}: $W(S) \in [0, W_{\max}]$ for all $S$.
  \item \textbf{Hysteresis stability}: No oscillation for fluctuations below $\delta$.
  \item \textbf{NaN propagation}: NaN in any dimension produces a valid (non-NaN) result.
\end{enumerate}

\subsection{Integration Tests}

450+ tests in the bridge-common crate cover the full surface area, including
cross-cluster dispatch, rate limiting, reputation operations, and credential
lifecycle.

\section{Comparison with Existing Approaches}

\begin{table*}[t]
\centering
\begin{tabular}{lcccc}
\toprule
\textbf{System} & \textbf{Sybil Resist.} & \textbf{Non-Plutocratic} & \textbf{Multi-Axis} & \textbf{Progressive} \\
\midrule
Token voting & Low & No & No & No \\
Quadratic voting & Low & Partial & No & Yes \\
Conviction voting & Low & Yes & No & Yes \\
Reputation-only & Medium & Yes & No & Yes \\
BrightID (social graph) & High & Yes & No & No \\
\textbf{Consciousness gating} & \textbf{High} & \textbf{Yes} & \textbf{Yes} & \textbf{Yes} \\
\bottomrule
\end{tabular}
\caption{Comparison of decentralized governance access control mechanisms.}
\label{tab:comparison}
\end{table*}

\section{Limitations and Future Work}

\begin{enumerate}
  \item \textbf{Weight selection}: The dimension weights (0.25/0.25/0.30/0.20)
    were determined empirically. Formal optimization using game-theoretic
    analysis of attack cost functions is future work.
  \item \textbf{Community dimension cold start}: New communities lack
    attestation data. The bootstrap mode mitigates but does not eliminate
    this problem.
  \item \textbf{Cross-community portability}: Consciousness profiles are
    community-specific. A mechanism for portable profiles across independent
    Mycelix deployments remains open.
  \item \textbf{Formal verification}: The Rust implementation has 450+ tests
    but no formal verification. Proving the security properties in a theorem
    prover (e.g., Coq or Lean) is desirable for high-assurance deployments.
\end{enumerate}

\section{Conclusion}

Consciousness gating provides a principled framework for decentralized governance
access control that is simultaneously Sybil-resistant, non-plutocratic,
multi-axis secure, and operationally efficient. The 4-dimensional profile
captures the essential qualities of a governance participant---verified identity,
demonstrated reputation, community recognition, and active engagement---while
the sigmoid-smoothed vote weighting and hysteresis mechanisms ensure stable,
continuous governance participation.

The implementation in Mycelix demonstrates that this framework is practical at
scale: 141+ zomes across 16 clusters, all sharing the same consciousness
gating infrastructure through a single bridge-common crate with 450+ tests.

\begin{thebibliography}{99}
  \bibitem{buterin2019} Buterin, V. (2019). Moving beyond coin voting governance.
    Blog post.
  \bibitem{weyl2018} Posner, E. A. and Weyl, E. G. (2018). \textit{Radical Markets}.
    Princeton University Press.
  \bibitem{conviction2019} Emmett, J. (2019). Conviction Voting: A Novel
    Continuous Decision Making Alternative. Commons Stack blog.
  \bibitem{ostrom1990} Ostrom, E. (1990). \textit{Governing the Commons}.
    Cambridge University Press.
  \bibitem{brock2018} Brock, A. and Harris-Braun, E. (2018). Holochain:
    Scalable Agent-Centric Distributed Computing.
  \bibitem{w3cdid} W3C (2022). Decentralized Identifiers (DIDs) v1.0.
  \bibitem{douceur2002} Douceur, J. R. (2002). The Sybil Attack. IPTPS.
  \bibitem{bright2020} BrightID (2020). Universal Proof of Uniqueness. Whitepaper.
\end{thebibliography}

\end{document}
